% !TeX root=../main.tex

\chapter{مقدمه}
\label{ch:intro}

\section{پیش‌گفتار}
مدل‌های زبانی\LTRfootnote{Language Model}
بزرگ در سال‌های اخیر به ابزاری اساسی برای طیف وسیعی از کاربردهای پردازش زبان طبیعی\LTRfootnote{Natural Language Processing}
تبدیل شده‌اند. این مدل‌ها که عموماً بر پایهٔ معماری ترنسفورمر\LTRfootnote{Transformer}~\cite{vaswani2017attention}
ساخته می‌شوند، با پیش‌آموزش\LTRfootnote{Pretraining}
بر روی حجم عظیمی از متن، بازنمایی‌های\LTRfootnote{Representations}
برداری‌ای از زبان می‌آموزند که در لایه‌های میانی و در جریان به‌اصطلاح باقیمانده\LTRfootnote{Residual Stream}
شبکه جریان می‌یابند. سنجهٔ نهایی موفقیت این مدل‌ها اغلب توانایی آن‌ها در پیش‌بینی نشانهٔ\LTRfootnote{Token}
بعدی متن است، اما آنچه در عمل کاربرد این مدل‌ها را در وظایفی مانند بازیابی اطلاعات، پرسش‌وپاسخ، خلاصه‌سازی و تشخیص هم‌ارزی معنایی\LTRfootnote{Semantic Equivalence}
تعیین می‌کند، کیفیت بازنمایی درونی‌ای است که این مدل‌ها از معنای جملات می‌سازند؛ به بیان دیگر، اینکه آیا فاصلهٔ میان بازنمایی دو جمله در فضای برداری مدل، واقعاً بازتابی از فاصلهٔ معنایی آن دو جمله هست یا صرفاً بازتابی از شباهت سطحی و واژگانی آن‌هاست.

پژوهش‌های متعددی نشان داده‌اند که این دو نوع شباهت--شباهت واژگانی و شباهت معنایی--می‌توانند به‌شدت از یکدیگر جدا باشند. مجموعه‌دادهٔ
\lr{PAWS}\LTRfootnote{Paraphrase Adversaries from Word Scrambling}~\cite{zhang2019paws}،
که مبنای عملی این پایان‌نامه نیز هست، دقیقاً برای آشکارسازی همین شکاف طراحی شده است: جفت‌جمله‌هایی با همپوشانی واژگانی بسیار بالا که با جابه‌جایی ترتیب کلمات یا ترجمهٔ رفت‌وبرگشتی\LTRfootnote{Back-translation}
ساخته شده‌اند، به‌گونه‌ای که برخی از این جفت‌ها هم‌معنی\LTRfootnote{Paraphrase}
و برخی دیگر، على‌رغم شباهت ظاهری بسیار زیاد، ناهم‌معنی هستند. برای نمونه، دو جمله که تنها در ترتیب قرارگیری یک عبارت مکانی با یکدیگر تفاوت دارند می‌توانند معنای کاملاً متفاوتی داشته باشند، در حالی‌که از منظر یک مدل متکی بر همپوشانی واژگانی، تقریباً یکسان به نظر می‌رسند. سازندگان
\lr{PAWS}
نشان دادند مدل‌های زمان خود که بر مجموعه‌داده‌های متعارف تشخیص بازنویسی\LTRfootnote{Paraphrase Identification}
آموزش دیده بودند، در تشخیص این جفت‌های دشوار به‌شدت ضعیف عمل می‌کنند (دقتی کمتر از ۴۰ درصد)~\cite{zhang2019paws}، و این ضعف نشانه‌ای از آن است که بازنمایی‌های آموخته‌شده توسط این مدل‌ها، حتی وقتی روی وظایف شباهت معنایی به ظاهر موفق عمل می‌کنند، ممکن است بیش از آنکه معنا را رمزگذاری کرده باشند، به الگوهای سطحی واژگانی متکی باشند.

نقطهٔ آغاز این پایان‌نامه، مواجهه با همین مجموعه‌داده و همین مشاهده بود: وقتی جفت‌جمله‌های
\lr{PAWS}
از یک مدل زبانی پیش‌آموزش‌دیده عبور داده شدند و بازنمایی آن‌ها در لایه‌های مختلف مدل بررسی شد، فاصلهٔ میان بازنمایی جفت‌های هم‌معنی و جفت‌های ناهم‌معنی، آن‌قدر که انتظار می‌رفت، از هم متمایز نبود. این مشاهده این پرسش را پیش کشید که آیا می‌توان مدل زبانی را طوری ریزتنظیم\LTRfootnote{Fine-tuning}
کرد که بازنمایی جفت‌جمله‌های هم‌معنی را به یکدیگر نزدیک‌تر و بازنمایی جفت‌جمله‌های ناهم‌معنی را از یکدیگر دورتر کند، بدون آنکه توانایی عمومی مدل در زبان‌مدل‌سازی\LTRfootnote{Language Modeling}
را به‌طور محسوسی از دست بدهد.

پاسخ ساده‌ای که برای این پرسش می‌توان متصور شد، اعمال مستقیم یک تابع هزینهٔ تقابلی\LTRfootnote{Contrastive Loss}
بر روی فضای پنهان\LTRfootnote{Latent Space}
خام خود مدل--یعنی همان جریان باقیمانده--است. اما فضای خام مدل، فضایی است که در آن هزاران مفهوم به‌طور همزمان و در قالب پدیده‌ای موسوم به برهم‌نهی\LTRfootnote{Superposition}~\cite{elhage2022superposition}
در ابعاد مشترکی جای گرفته‌اند؛ در نتیجه، فشردن یا کشیدن این فضا در راستای یک هدف مشخص، می‌تواند به شکلی غیرقابل پیش‌بینی، مفاهیم نامرتبط دیگری را نیز جابه‌جا کند. در سال‌های اخیر، خودرمزگذارهای پراکنده\LTRfootnote{Sparse Autoencoders, SAE}
به‌عنوان ابزاری برای گشودن این برهم‌نهی و استخراج مؤلفه‌هایی نسبتاً تک‌معنا\LTRfootnote{Monosemantic}
از فضای فعال‌سازی مدل‌های زبانی معرفی شده‌اند~\cite{bricken2023monosemanticity,cunningham2023sae}، و تاکنون عمدتاً برای اهداف تحلیلی و تفسیرپذیری\LTRfootnote{Interpretability}
مورد استفاده قرار گرفته‌اند: یعنی برای فهمیدن آنچه یک مدل از پیش آموزش‌دیده «می‌داند»، نه برای تغییردادن آنچه مدل می‌آموزد.

ایدهٔ محوری این پایان‌نامه، جابه‌جایی نقش خودرمزگذار پراکنده از یک ابزار صرفاً تحلیلی به یک فضای هدف برای آموزش است: به‌جای اعمال تابع هزینهٔ ناوردایی\LTRfootnote{Invariance Loss}
مستقیماً بر فضای خام مدل، این تابع هزینه بر روی کدهای نهان حاصل از عبوردادن فعال‌سازی‌های مدل از یک خودرمزگذار پراکندهٔ از پیش‌آموزش‌دیده و ثابت\LTRfootnote{Frozen}
اعمال می‌شود. از آنجا که این خودرمزگذار، بازنمایی خام را به مجموعه‌ای پراکنده‌تر و به‌طور نسبی مجزاتر از ویژگی‌ها\LTRfootnote{Features}
تجزیه می‌کند، این پایان‌نامه این فرضیه را می‌آزماید که نزدیک‌کردن جفت‌جمله‌های هم‌معنی به هم و دورکردن جفت‌جمله‌های ناهم‌معنی از هم، در این فضای پراکنده، مؤثرتر و «تمیزتر» از انجام همان کار در فضای خام مدل باشد.

\section{بیان مسئله}
مسئلهٔ این پژوهش را می‌توان در سه پرسش پژوهشی صورت‌بندی کرد:

\begin{enumerate}
\item آیا اعمال یک تابع هزینهٔ ناوردایی--که جفت‌جمله‌های هم‌معنی
\lr{PAWS}
را نزدیک و جفت‌جمله‌های ناهم‌معنی را دور از هم نگاه می‌دارد--در فضای نهان یک خودرمزگذار پراکندهٔ ثابت، نسبت به اعمال همان تابع هزینه در فضای خام جریان باقیمانده، جداسازی بهتری میان جفت‌های هم‌معنی و ناهم‌معنی تولید می‌کند؟
\item این جداسازی چگونه به انتخاب لایه(های) هدف در عمق مدل و به وزن نسبی تابع هزینهٔ ناوردایی در برابر تابع هزینهٔ زبانی وابسته است؟
\item ریزتنظیم‌کردن مدل در جهت این هدف، با چه هزینه‌ای--از نظر افت کیفیت زبان‌مدل‌سازی و از نظر اعتبار خود خودرمزگذار--همراه است، و آیا می‌توان این هزینه را با یک منظم‌ساز صریح مهار کرد؟
\end{enumerate}

پرسش سوم از این جهت اهمیت دارد که کل رویکرد پیشنهادی به یک فرض ضمنی متکی است: اینکه خودرمزگذار پراکندهٔ ثابت، در طول ریزتنظیم مدل زبانی، همچنان بازسازی‌کنندهٔ باکیفیتی از فعال‌سازی‌های مدل باقی بماند. اگر فعال‌سازی‌های مدل، در پی ریزتنظیم، از ناحیه‌ای که خودرمزگذار روی آن آموزش دیده دور شوند، سنجه‌هایی که در فضای این خودرمزگذار محاسبه می‌شوند (از جمله همان سنجهٔ جداسازی هم‌معنی/ناهم‌معنی که هدف اصلی این پژوهش است) دیگر قابل اعتماد نخواهند بود، چرا که آنچه اندازه‌گیری می‌شود، بازتابی از یک ابزار سنجش نامنطبق خواهد بود، نه بازتابی از تغییر واقعی در بازنمایی مدل. به همین دلیل، این پژوهش از ابتدا یک جزء منظم‌ساز بازسازی\LTRfootnote{Reconstruction Regularizer}
را نیز در تابع هزینهٔ کلی گنجانده است تا اعتبار خودرمزگذار، در طول آموزش و پایش حفظ شود.

\section{هدف تحقیق}
هدف اصلی این پایان‌نامه، طراحی و ارزیابی تجربی روشی برای ریزتنظیم کارآمد پارامتر\LTRfootnote{Parameter-Efficient Fine-Tuning}
یک مدل زبانی از پیش‌آموزش‌دیده است، به‌گونه‌ای که بازنمایی جفت‌جمله‌های هم‌معنی و ناهم‌معنی--آن‌گونه که در مجموعه‌دادهٔ
\lr{PAWS}
تعریف شده‌اند--در فضای نهان یک خودرمزگذار پراکندهٔ ثابت، از یکدیگر بهتر متمایز شوند. این هدف کلی به اهداف فرعی زیر تجزیه می‌شود:

\begin{itemize}
\item طراحی یک تابع هزینهٔ ترکیبی که هزینهٔ زبان‌مدل‌سازی روی یک پیکرهٔ عمومی (برای حفظ سلامت زبانی مدل)، هزینهٔ ناوردایی روی جفت‌جمله‌های
\lr{PAWS}
(برای شکل‌دادن به بازنمایی)، و منظم‌ساز بازسازی خودرمزگذار (برای حفظ اعتبار) را همزمان بهینه کند.
\item مقایسهٔ تجربی اعمال تابع هزینهٔ ناوردایی در فضای خام در برابر فضای نهان خودرمزگذار، تحت شرایط کنترل‌شدهٔ یکسان.
\item بررسی حساسیت نتایج به دو ابرپارامتر کلیدی: لایهٔ (یا لایه‌های) هدف اعمال تابع هزینهٔ ناوردایی، و وزن نسبی این تابع هزینه.
\item تحلیل اثرات جانبی این ریزتنظیم--از جمله میزان تعمیم‌پذیری بهبود به لایه‌های غیرهدف، تأثیر آن بر سرگشتگی\LTRfootnote{Perplexity}
مدل روی متن عمومی، و پایداری کیفیت بازسازی خودرمزگذار در طول آموزش.
\end{itemize}

\section{روش انجام تحقیق (خلاصه)}
این پژوهش از نوع تجربی\LTRfootnote{Empirical}
و کاربردی است. مدل زبانی
\lr{Pythia\nobreakdash-70m\nobreakdash-deduped}~\cite{biderman2023pythia}
به همراه یک تطبیق‌دهندهٔ کم‌رتبه\LTRfootnote{LoRA, Low-Rank Adaptation}~\cite{hu2021lora}
به‌عنوان مدل پایه انتخاب شد. برای هر لایه از این مدل، از یک خودرمزگذار پراکندهٔ از پیش‌آموزش‌دیده و منتشرشدهٔ عمومی--بارگذاری‌شده از طریق کتابخانهٔ
\lr{SAELens}~\cite{bloom2024saelens}--استفاده شده که در طول کل فرایند ریزتنظیم، ثابت (منجمد) نگه داشته می‌شود. جفت‌جمله‌های هم‌معنی و ناهم‌معنی مجموعه‌دادهٔ
\lr{PAWS}~\cite{zhang2019paws}
منبع سیگنال نظارتی برای تابع هزینهٔ ناوردایی هستند، و بخشی از پیکرهٔ
\lr{WikiText-2}
برای حفظ سلامت زبان‌مدل‌سازی مدل در کنار آن به کار می‌رود. چهار شرط آزمایشی متفاوت (از حالتی که هیچ تابع هزینهٔ ناوردایی‌ای اعمال نمی‌شود، تا حالت‌هایی که این تابع هزینه در فضای خام یا فضای خودرمزگذار، با یا بدون یک جزء تکمیلی حمایتی، اعمال می‌شود) در برابر پنج مقدار برای وزن تابع هزینهٔ ناوردایی و ده حالت متفاوت برای انتخاب لایهٔ هدف (شامل تک‌لایه‌ها و ترکیب چند لایهٔ همزمان) به‌صورت آزمایش‌های شبکه‌ای\LTRfootnote{Grid Search}
سیستماتیک اجرا و ارزیابی شده‌اند. جزئیات کامل این روش، موضوع فصل سوم است.

\section{نوآوری و اهمیت پژوهش}
تا آنجا که در چارچوب این پایان‌نامه بررسی شده، بیشتر کاربردهای موجود خودرمزگذارهای پراکنده در ادبیات تفسیرپذیری مکانیستی\LTRfootnote{Mechanistic Interpretability}، یا برای \emph{کشف} و نام‌گذاری پس‌ازآموزش ویژگی‌های یک مدل ثابت به کار رفته‌اند~\cite{bricken2023monosemanticity,templeton2024scaling}، یا برای \emph{هدایت}\LTRfootnote{Steering}
رفتار مدل در زمان استنتاج\LTRfootnote{Inference}
با دخالت مستقیم در مقدار یک ویژگی از پیش‌شناسایی‌شده~\cite{turner2023actadd,templeton2024scaling}. نوآوری این پایان‌نامه در آن است که خودرمزگذار پراکنده را نه به‌عنوان ابزار تحلیل پس‌ازآموزش و نه به‌عنوان اهرم دستکاری زمان استنتاج، بلکه به‌عنوان \emph{فضای هدف یک تابع هزینه در حین آموزش} به کار می‌گیرد؛ به بیان دیگر، خودرمزگذار در این طرح، نقش یک عدسی ثابت را ایفا می‌کند که از میان آن، سیگنال گرادیانی حاصل از یک وظیفهٔ نظارت‌شده (تشخیص هم‌معنی‌بودن) به پارامترهای مدل زبانی بازمی‌گردد.

سهم دیگر این پژوهش، ارائهٔ یک منظم‌ساز صریح برای حفظ اعتبار خودرمزگذار در حین این نوع آموزش است؛ نکته‌ای که در کاربردهای رایج (کشف پس‌ازآموزش) اساساً محل بحث نیست، چرا که در آن‌ها مدل زبانی حین تحلیل تغییر نمی‌کند. سرانجام، این پژوهش با انجام یک مطالعهٔ حذفی\LTRfootnote{Ablation Study}
نسبتاً گسترده--در دو بُعد عمق لایه و وزن تابع هزینه، و در برابر یک شرط کنترلی فاقد هرگونه تابع هزینهٔ ناوردایی--تصویری تجربی از سازوکارها، محدودیت‌ها و مصالحه‌های\LTRfootnote{Trade-offs}
این رویکرد به دست می‌دهد که می‌تواند مبنایی برای پژوهش‌های آتی روی مدل‌ها و مجموعه‌داده‌های بزرگ‌تر باشد.

\section{تعریف واژه‌های کلیدی}
برای روشن‌شدن ادامهٔ متن، برخی از مهم‌ترین اصطلاحاتی که در این پایان‌نامه به‌کرّات استفاده می‌شوند، در این‌جا به‌اختصار معرفی می‌شوند؛ تعریف کامل‌تر و فنی‌تر هر یک، در فصل‌های دوم و سوم ارائه خواهد شد.

\begin{description}
\item[بازنویسی\LTRfootnote{Paraphrase}:] دو جمله که با وجود تفاوت در واژگان یا ساختار سطحی، معنای یکسان یا تقریباً یکسانی را منتقل می‌کنند.
\item[جریان باقیمانده\LTRfootnote{Residual Stream}:] مسیر برداری اصلی‌ای که در معماری ترنسفورمر از ورودی تا خروجی شبکه امتداد دارد و هر لایه، خروجی خود را به‌صورت جمعی به آن می‌افزاید.
\item[خودرمزگذار پراکنده (\lr{SAE})\LTRfootnote{Sparse Autoencoder}:] شبکه‌ای کوچک که یک بردار فعال‌سازی را به فضایی با ابعاد بسیار بیشتر ولی با فعالیت پراکنده (اکثر مؤلفه‌ها صفر) می‌برد و سپس تلاش می‌کند از روی همان کد پراکنده، بردار اصلی را بازسازی کند؛ در این پایان‌نامه، این شبکه از پیش آموزش دیده و در حین ریزتنظیم مدل زبانی، ثابت نگه داشته می‌شود.
\item[تابع هزینهٔ ناوردایی\LTRfootnote{Invariance Loss}:] جزئی از تابع هزینهٔ کلی که بازنمایی جفت‌جمله‌های هم‌معنی را به هم نزدیک و بازنمایی جفت‌جمله‌های ناهم‌معنی را از هم دور نگاه می‌دارد.
\item[لورا (\lr{LoRA})\LTRfootnote{Low-Rank Adaptation}:] روشی برای ریزتنظیم کارآمد پارامتر که به‌جای به‌روزرسانی مستقیم کل ماتریس‌های وزن مدل، یک به‌روزرسانی کم‌رتبه به آن‌ها می‌افزاید و تنها همین بخش کوچک را آموزش می‌دهد.
\item[\lr{AUROC}:] سطح زیر منحنی مشخصهٔ عملکرد گیرنده\LTRfootnote{Receiver Operating Characteristic}، سنجه‌ای بین صفر تا یک برای کیفیت تفکیک دو دسته (در اینجا: جفت‌های هم‌معنی در برابر ناهم‌معنی) بر اساس یک سنجهٔ پیوسته (در اینجا: فاصلهٔ کسینوسی)؛ مقدار ۵٫۰ برابر با حدس تصادفی و مقدار ۱ برابر با تفکیک کامل است~\cite{fawcett2006roc}.
\end{description}

\section{ساختار پایان‌نامه}
باقی این پایان‌نامه در چهار فصل سازمان‌دهی شده است. فصل دوم مروری بر ادبیات موضوع را در چهار محور اصلی ارائه می‌دهد: بازنمایی و شباهت معنایی در مدل‌های زبانی، یادگیری تقابلی، تفسیرپذیری مکانیستی و خودرمزگذارهای پراکنده، و ریزتنظیم کارآمد پارامتر و حفظ رفتار مدل در حین آن. فصل سوم روش پیشنهادی--معماری، تابع هزینه، شرایط آزمایشی و طرح آزمایش‌های شبکه‌ای--را با جزئیات کامل شرح می‌دهد. فصل چهارم نتایج تجربی حاصل از این آزمایش‌ها را، به همراه تحلیل و نمودارهای مرتبط، ارائه می‌کند. فصل پنجم پایان‌نامه را با جمع‌بندی یافته‌ها، بیان محدودیت‌ها و پیشنهادهایی برای پژوهش‌های آینده به پایان می‌برد.
